\subsection{深度学习 v.s.传统神经网络}

早期的神经网络由于计算资源的限制和优化方法的不成熟，难以扩展到实际应用场景。特别是在深度方面的局限，使得这些模型在复杂任务上表现不佳。直到深度学习的出现，神经网络才真正突破瓶颈，在多个领域取得革命性进展。那么，深度学习与传统神经网络的关键区别是什么？


%--
首先，深度学习突破了浅层神经网络的深度限制，提高了模型的表达能力。 早期的神经网络，如多层感知机（MLP），通常只包含一两个隐藏层，受限于计算能力和优化方法，训练更深的网络变得极为困难。然而，研究者们发现，单隐藏层的神经网络虽然理论上可以逼近任何连续函数，但其参数数量需要指数级增长，导致计算成本极高，难以扩展到复杂任务。相比之下，深度学习通过增加网络层数，使得模型能够更高效地学习复杂模式，从而在语音、图像、文本等领域展现出远超浅层网络的能力
%--

首先，{\bf 传统神经网络的局限性在于网络深度有限，依赖手工特征工程。} 
此外，这些方法依赖人工特征提取，例如在图像处理任务中，需要研究者手动设计边缘、角点、纹理等特征，再输入神经网络进行分类。这种方法不仅依赖领域专家的经验，还难以适应复杂多变的任务场景。

相比之下，{\bf 深度学习采用多层神经网络结构，能够端到端学习数据的层级特征}，避免了人工特征工程的瓶颈。以卷积神经网络（CNN）为例，底层神经元可以自动学习简单的边缘特征，中层神经元学习局部纹理，高层神经元则能捕捉整体形状和语义信息。这样的层次化特征学习方式，使得深度学习在高维数据（如图像、语音、文本）处理中表现出极强的能力。此外，深度学习可以通过自监督学习和无监督学习的方式，在海量数据中发现潜在模式，使得模型具有更强的泛化能力。

%--
其次，深度学习摆脱了传统神经网络对人工特征工程的依赖。 在传统机器学习方法中，研究者需要人为设计特征，例如在图像处理中，需要手动提取边缘、颜色直方图、纹理等特征，再输入神经网络进行分类。 这种方法不仅依赖领域知识，还难以适应不同类型的数据。
深度学习通过端到端的训练方式，使得神经网络可以直接从原始数据中自动提取特征。例如，在卷积神经网络（CNN）中：
- 底层神经元 负责检测边缘、角点等基本特征；
- 中层神经元 学习更复杂的局部模式，如纹理、形状；
- 高层神经元 则学习全局语义信息，如物体类别、人脸等。

这种层次化的特征表示，使得深度学习不仅能适应不同任务，而且能在高维复杂数据（如图像、语音、文本）处理中展现出极强的能力。
%--

另一个关键区别在于{\bf 优化能力的提升}。早期神经网络受限于梯度消失、局部最优等问题，而深度学习的优化方法更为先进，如残差连接（ResNet）和自适应梯度算法，使得训练更深层的网络成为可能。这不仅提升了模型的学习能力，还扩展了神经网络在各类任务上的适用范围。

%--
最后，优化方法的突破是深度学习成功的关键。 传统神经网络的训练往往面临梯度消失、局部最优等问题，而深度学习的兴起伴随着一系列优化方法的突破：
- ReLU（修正线性单元）激活函数：缓解梯度消失问题，使得深层网络训练更加稳定。
- 批量归一化（Batch Normalization）：加速训练并稳定梯度，使得深度模型更容易收敛。
- 残差网络（ResNet）：引入“跳跃连接”（Skip Connections），使得梯度可以顺利传播到前层，从而成功训练上百层的深度神经网络。
%--

这些技术的突破，使得深度学习不仅能够训练更深的网络，还显著提高了训练效率和模型的稳定性，进一步推动了人工智能的发展。
